
\subsubsection{Experimental Questions}

Our experimental design operationalizes the semantic validity hypothesis through three interrelated questions:

\textbf{Q1: Does horizontal flip cause differential degradation?} We predict asymmetric digit accuracy will decrease significantly under flip augmentation compared to baseline, while symmetric digit accuracy remains stable. We test this via per-class accuracy grouped by digit symmetry at flip probability p=0.5.

\textbf{Q2: Is degradation dose-dependent?} If label noise is the mechanism, degradation magnitude should be a monotonic function of flip probability. We test flip probabilities p $\in$ {0.3, 0.5, 0.9} plus baseline (p=0.0), computing Spearman rank correlation.

\textbf{Q3: Is the effect specific to semantic invalidity?} To isolate semantic invalidity from general augmentation effects, we compare horizontal flip against rotation $\pm 15$°. If semantic invalidity is the causal factor, rotation should show no differential effect on asymmetric versus symmetric digits.

\subsubsection{Experimental Protocol}

All experiments use MNIST (60,000 train / 10,000 test images, $28\times 28$ grayscale, 10 classes). Model architecture is a standard convolutional neural network (PyTorch MNISTNet): two convolutional layers ($32\rightarrow 64$ channels), ReLU activations, $2\times 2$ max pooling, dropout (0.25 after conv, 0.5 after first FC layer), and two fully connected layers ($128\rightarrow 10)$. This ~100K-parameter architecture is deliberately shallow to test the effect on resource-constrained models.

Training configuration is maintained within sub-hypotheses. Primary experiments (h-e1, h-m) use Adadelta optimizer (lr=1.0), StepLR scheduler (step=1, $\gamma =0.7)$, batch size 64, 14 epochs, NLLLoss. Mechanism validation (h-m1) uses Adam optimizer (lr=0.001), StepLR scheduler ($\gamma =0.7)$, batch size 64, early stopping (patience=5), NLLLoss.

Augmentation conditions tested:
\begin{itemize}
\item \textbf{Baseline}: ToTensor + Normalize only
\item \textbf{Flip30, Flip50, Flip90}: Baseline + RandomHorizontalFlip(p={0.3, 0.5, 0.9})
\item \textbf{Rotation}: Baseline + RandomRotation(degrees=15)
\end{itemize}

Evaluation metrics: per-class test accuracy for all 10 digits, then aggregate by symmetry groups: asymmetric = mean({2,3,5,6,7,9}), symmetric = mean({0,1,8}). For differential effect, we use Wilcoxon signed-rank test with significance threshold p<0.05. For dose-response, we compute Spearman rank correlation across conditions $\times$ seeds, testing null hypothesis $\rho$=0 against alternative $\rho$<0.

\subsubsection{Controlled Design Rationale}

\textbf{Positive Control (Rotation)}: Rotation $\pm 15$° serves as a ''semantically valid augmentation'' control. Unlike flip, rotation preserves digit class identity. If degradation were due to general augmentation effects, we should see similar degradation under rotation. Observing neutral rotation effects while flip causes degradation confirms semantic invalidity as the causal factor.

\textbf{Negative Control (Symmetric Digits)}: Symmetric digits {0,1,8} are flip-invariant by geometry. Under our label noise mechanism, these digits should show minimal degradation because flip does not create semantically invalid images for symmetric classes. This within-dataset control strengthens causal inference.

\textbf{Dose-Response Design}: Testing multiple flip probabilities {0.3, 0.5, 0.9} rather than a single binary comparison provides mechanistic evidence. If label noise is the causal mechanism, degradation should scale with flip probability. A monotonic dose-response relationship is stronger evidence for a causal mechanism than a single effect size measurement.

